\newcommand{\nomeEstado}{BAHIA}
\newcommand{\siglaEstado}{BA}
\newcommand{\nomeConcessionaria}{NEOENERGIA COELBA}
\newcommand{\cnpjConcessionaria}{15.139.629/0001-94}
\newcommand{\presidenteConcessionaria}{Sr. Thiago Guth.}

\newcommand{\empresaResponsavel}{ABEL}
\newcommand{\nomeMunicipio}{Santo Antônio de Jesus}
\newcommand{\nomeMunicipioMaisculo}{\MakeUppercase{\nomeMunicipio}}
\newcommand{\cnpjMunicipio}{13.825.476/0001-03}
\newcommand{\sedePrefeitura}{Avenida Ursicino Pinto de Queiroz, nº 167, Centro}
\newcommand{\generoPrefeitura}{pelo Prefeito}
\newcommand{\nomePrefeito}{Genival Deolino Souza}
\newcommand{\numeroContrato}{248/2024} % Número do contrato (Não precisa por "Nº")
\newcommand{\quantidadeAditivos}{4} % Quantidade de aditivos do contrato
\newcommand{\legitimidade}{3} % 0 - Sem legitimidade, C - No contrato, 1,2,3... - No Aditivo (colocar o número do aditivo)
\newcommand{\publicacao}{Diário Oficial do Município}
\newcommand{\dataPublicacao}{14 de maio de 2024} % Data da publicação
\newcommand{\tituloClausula}{CLÁUSULA PRIMEIRA – DO OBJETO } % Título da cláusula que confere poderes de representação
\newcommand{\clausula}{CLÁUSULA PRIMEIRA – DO OBJETO 

O presente Termo Aditivo refere-se ao ajuste na redação do objeto do Contrato nº 248/2024, celebrado com a empresa ABEL CUNHA – SOCIEDADE INDIVIDUAL DE ADVOCACIA, visando adequá-lo às necessidades técnicas identificadas na fase de execução contratual, diante da constatação de que o contrato original não contempla, de forma expressa, a possibilidade de representação administrativa do Município de Santo Antônio de Jesus – BA perante a Agência Nacional de Energia Elétrica – ANEEL. A alteração proposta tem como objetivo garantir a previsão contratual dessa atuação, em conformidade com os procedimentos e prazos estabelecidos na Resolução Normativa ANEEL nº 1.000/2021, conforme justificativa apresentada no Processo Administrativo nº 6.701/2025.} % Texto da cláusula que confere poderes de representação

\newcommand{\telefone}{7536321330}
\newcommand{\email}{falecom@saj.ba.gov.br}